% =====================================================================
%  CARNET D'ENTRAÎNEMENT — CHIMIE — FICHE 3 — THÈME 1
%  Particules subatomiques — NIVEAU DIFFICILE — ÉNONCÉS
% =====================================================================
\documentclass[11pt,a4paper]{article}

\usepackage[utf8]{inputenc}
\usepackage[T1]{fontenc}
\usepackage{lmodern}
\usepackage[french]{babel}

\usepackage{amsmath,amssymb,mathtools}
\usepackage{xcolor}
\usepackage{colortbl}
\usepackage{tikz}
\usetikzlibrary{arrows.meta,positioning,calc}
\usepackage[most]{tcolorbox}
\usepackage{enumitem}
\usepackage{booktabs}
\usepackage{array}
\usepackage[a4paper,margin=2.1cm,headheight=26pt,headsep=9pt]{geometry}
\usepackage{fancyhdr}
\usepackage[hidelinks]{hyperref}

\definecolor{primary}{RGB}{124,78,140}
\definecolor{primarydark}{RGB}{74,44,92}
\definecolor{defcol}{RGB}{0,114,178}
\definecolor{thmcol}{RGB}{0,140,120}
\definecolor{attncol}{RGB}{200,40,55}

\newcommand{\nuc}[3]{\prescript{#1}{#2}{\mathrm{#3}}}
\newcommand{\pp}{p^{+}}
\newcommand{\ee}{e^{-}}
\newcommand{\nz}{n^{0}}

\pagestyle{fancy}
\fancyhf{}
\fancyhead[L]{\small\textcolor{primarydark}{\textbf{Fiche 3 — Thème 1}} : Particules subatomiques}
\fancyhead[R]{\small\textcolor{primarydark}{\textsc{Difficile — Énoncés}}}
\fancyfoot[C]{\small\thepage}
\renewcommand{\headrulewidth}{0.8pt}
\renewcommand{\headrule}{\hbox to\headwidth{\color{primary}\leaders\hrule height \headrulewidth\hfill}}

\newcounter{exo}
\newcommand{\exo}[1][]{\stepcounter{exo}\par\medskip%
  \noindent{\color{primary}\rule[-2pt]{3.5pt}{15pt}}\ \textbf{\color{primarydark}\large Exercice \theexo}%
  \ifx\relax#1\relax\relax\else\ \textmd{\normalsize\color{black!65}— \itshape #1}\fi\par\nopagebreak\smallskip}

\setlist{leftmargin=1.7em,topsep=2pt,itemsep=3pt}

\begin{document}

\begin{tikzpicture}
\node[rectangle,rounded corners=6pt,fill=primary,text=white,
      minimum width=\textwidth,minimum height=1.35cm,align=center,inner sep=6pt]
  {\Large\bfseries Thème 1 — Particules subatomiques\\[2pt]
   \normalsize\mdseries Niveau Difficile — Énoncés};
\end{tikzpicture}

\vspace{3pt}
{\small\itshape Données : $e = 1{,}6\times10^{-19}$ C ; $m(\ee) = 9{,}11\times10^{-31}$ kg ;
$m(\pp)\approx 1{,}7\times10^{-27}$ kg ; $1\ \text{eV} = 1{,}66\times10^{-19}$ J ;
vitesse de la lumière $c \approx 3{,}0\times10^{8}$ m/s. Les sous-questions se construisent : traite-les
dans l'ordre.}

% ---------------------------------------------------------------
\exo[vitesse d'un électron à partir de son énergie cinétique]
Un électron possède une énergie cinétique $E_c = 18{,}22$ eV. On rappelle $E_c = \tfrac12 m v^2$.
\begin{enumerate}[label=\alph*)]
  \item Convertis $E_c$ en joules.
  \item Isole la vitesse $v$ dans la relation $E_c = \tfrac12 m v^2$, puis remplace par les valeurs
        (sans encore calculer le résultat final).
  \item Calcule $v$. \emph{Indication} : simplifie d'abord le facteur $\dfrac{2\times18{,}22}{9{,}11}$
        et le rapport des puissances de $10$.
  \item À quelle fraction (en \%) de la vitesse de la lumière $c$ cette vitesse correspond-elle ?
\end{enumerate}

% ---------------------------------------------------------------
\exo[le chemin inverse : de la vitesse à l'énergie]
Un autre électron se déplace à la vitesse $v = 3{,}0\times10^{6}$ m/s.
\begin{enumerate}[label=\alph*)]
  \item Calcule son énergie cinétique $E_c$, en joules.
  \item Exprime cette énergie en électron-volts (eV).
\end{enumerate}

% ---------------------------------------------------------------
\exo[une microbille chargée]
Une microbille porte une charge électrique de $-1{,}6\times10^{-15}$ C. Cette charge provient d'un
déséquilibre entre ses protons et ses électrons.
\begin{enumerate}[label=\alph*)]
  \item S'agit-il d'un excès de \textbf{protons} ou d'\textbf{électrons} ? Justifie par le signe.
  \item À combien de \textbf{charges élémentaires} $e$ en excès cette charge correspond-elle ?
  \item Combien d'électrons (ou de protons) en excès la microbille porte-t-elle ?
\end{enumerate}

% ---------------------------------------------------------------
\exo[électron contre proton : même vitesse, quelle énergie ?]
Un électron et un proton se déplacent à la \textbf{même vitesse} $v$.
\begin{enumerate}[label=\alph*)]
  \item En utilisant $E_c = \tfrac12 m v^2$, exprime le rapport $\dfrac{E_c(\pp)}{E_c(\ee)}$ en fonction
        des masses seules.
  \item Calcule numériquement ce rapport. Laquelle des deux particules possède la plus grande énergie
        cinétique, et environ combien de fois ?
\end{enumerate}

\end{document}
